\section{Small meditations}

\subsection{The democracy of the unimaginable}
From a human standpoint a googol, a googolplex, level \(d\), level \(p\), the crypto coma,
and even Graham's number blur into one hazy category: ``roughly infinity''. Intuition has a
ceiling --- a cognitive horizon --- and it sits absurdly low, somewhere around a few
thousand. Beyond it every giant looks alike. The difference between them is vast, real, and
invisible to us.

\subsection{A name instead of a number}
The crypto coma cannot be held --- so we hold its name: \(\ccoma\), three strokes. The name
is finite; the number cannot be written. To name is a way to tame what cannot be carried.
But the leash is not the beast. What do we hold when we write \(\ccoma\) --- a number, or a
word standing where the number will not fit?

\subsection{Almost every giant is nameless}
Short descriptions are finite in number: there simply are not infinitely many phrases
shorter than a given length. Numbers, however, are infinite. So almost every number ---
including almost every number the size of \(d\) --- can be neither named, nor described, nor
singled out. The crypto coma is a rare named giant in an infinite crowd of anonymous ones.

\subsection{Physically transcendent}
A computer running until the heat death of the Universe, flipping a bit every Planck
instant, would not print the digits of level \(d\) --- not even the count of its digits.
Computation has a physical limit. So some finite numbers are real yet forever
unrealizable: true --- and never to be built, counted, or shown anywhere in this cosmos.

\subsection{A ladder without a top}
\(z\) is the last letter, not the last number. Apply \(\star\) once more --- and you have
passed the crypto coma; once more --- and it lies far below. Name any ceiling, and you have
just laid the floor of the next room. There is no largest number --- only the largest of
those you have bothered to name.

\subsection{Discovered or invented?}
Did the crypto coma wait in a Platonic heaven for two boys on a bench? Or did it flash into
being the moment it was written down? An ultrafinitist would say: numbers we cannot reach
simply do not exist. The boys cheerfully invented a nightmare for philosophy.
